\documentclass[11pt]{article}
\usepackage{mathblatt}
\begin{document}
\weit
\typeout{PROBE-BEGIN flaechen-zone-f6-v1}
\begin{aufgabe}{\detokenize{flaechen-zone-f6-v1}}
\begin{teile}
\teil Probe

\winkel{90}{\alpha}
\end{teile}
\end{aufgabe}
\typeout{PROBE-END flaechen-zone-f6-v1}
\typeout{PROBE-BEGIN kreis-zone-f5-v3}
\begin{aufgabe}{\detokenize{kreis-zone-f5-v3}}
\begin{teile}
\teil Probe

\winkel[4]{65}{\alpha}
\end{teile}
\end{aufgabe}
\typeout{PROBE-END kreis-zone-f5-v3}
\typeout{PROBE-BEGIN symmetrie-abbildungen-zone-f6-v1}
\begin{aufgabe}{\detokenize{symmetrie-abbildungen-zone-f6-v1}}
\begin{teile}
\teil Probe

\winkel{60}{\alpha}
\end{teile}
\end{aufgabe}
\typeout{PROBE-END symmetrie-abbildungen-zone-f6-v1}
\typeout{PROBE-BEGIN trigonometrie-zone-f2-v3}
\begin{aufgabe}{\detokenize{trigonometrie-zone-f2-v3}}
\begin{teile}
\teil Probe

\winkel{135}{\varepsilon}
\end{teile}
\end{aufgabe}
\typeout{PROBE-END trigonometrie-zone-f2-v3}
\typeout{PROBE-BEGIN trigonometrische-funktionen-zone-f5-v4}
\begin{aufgabe}{\detokenize{trigonometrische-funktionen-zone-f5-v4}}
\begin{teile}
\teil Probe

\winkel{110}{\alpha}
\end{teile}
\end{aufgabe}
\typeout{PROBE-END trigonometrische-funktionen-zone-f5-v4}
\end{document}
